% ============================================
% MEM-20
% Gerado de molambudos.md — build_miolo.py
% ============================================
\section*{\textbf{MEM-20} --- A Morte do Vizinho de Cela}
\markboth{MEM-20 --- A Morte do Vizinho de Cela}{}

\noindent \textit{O único amigo que tive}

\noindent \textbf{1963 --- O único amigo que tive}

\vspace{0.5em}\noindent\rule{\textwidth}{0.5pt}\vspace{0.5em}

\noindent Ele se chamava Severino.

\noindent Paraense. 52 anos. Internado em 1935 porque a família não queria mais cuidar dele. Diagnóstico: \textquotedblleft{}idiotia\textquotedblright{}. O diagnóstico real: gagueira e epilepsia leve. Mas nos anos 30, gagueira era tratada com eletrochoque. Depois de 30 choques, Severino não gaguejava mais. Nem falava. Nem mexia.

\noindent Ele ficava na cela ao lado da minha.

\noindent Durante 20 anos, a gente conversava batendo na parede.

\noindent Código:
1 batida = sim
2 batidas = não
3 batidas = como vai você
4 batidas = estou com fome
5 batidas = estou com medo

\noindent A gente desenvolveu uma linguagem. Batidas longas, batidas curtas. Dava para dizer frases complexas. \textit{\textquotedblleft{}O enfermeiro veio hoje\textquotedblright{}} --- 3 longas, 2 curtas, pausa, 4 curtas. Eu entendia.

\noindent Severino era a única pessoa no Colônia que me tratava como gente.

\noindent Ele não olhava para o meu olho amarelo com medo. Ele olhava como quem olha para um amigo.

\vspace{0.5em}\noindent\rule{\textwidth}{0.5pt}\vspace{0.5em}

\noindent Em 1963, Severino pegou pneumonia.

\noindent O Colônia não tratava pneumonia. Deixava morrer. Morriam 30 pacientes por mês de pneumonia, tuberculose, desnutrição. O Colônia chamava de \textquotedblleft{}seleção natural\textquotedblright{}.

\noindent Eu ouvi a tosse dele pela parede.

\noindent A tosse durou uma semana. Cada dia mais fraca. No sétimo dia, a tosse parou.

\noindent Bati na parede.

\noindent 1 batida. \textit{\textquotedblleft{}Severino?\textquotedblright{}}

\noindent Silêncio.

\noindent 2 batidas. \textit{\textquotedblleft{}Acorda.\textquotedblright{}}

\noindent Silêncio.

\noindent 3 batidas. \textit{\textquotedblleft{}Como vai você?\textquotedblright{}}

\noindent Silêncio.

\noindent Eu sabia que ele tinha morrido. Mas fiquei batendo. Bati a noite inteira. Diferentes códigos. Diferentes perguntas. \textit{\textquotedblleft{}Você está ai?\textquotedblright{}\textit{ }\textquotedblleft{}Responde.\textquotedblright{}\textit{ }\textquotedblleft{}Não me deixa sozinho.\textquotedblright{}}

\noindent Ninguém respondeu.

\noindent De manhã, os enfermeiros vieram buscar o corpo.

\noindent Eu ouvi a porta da cela abrir. Ouvi o arrastar dos pés no chão. Depois, silêncio.

\noindent A cela ao lado ficou vazia.

\noindent Por três meses, ninguém ocupou.

\noindent Eu batia na parede de vez em quando. Não para conversar --- para ouvir o som do oco.

\vspace{0.5em}\noindent\rule{\textwidth}{0.5pt}\vspace{0.5em}

\noindent Depois que Severino morreu, a criatura falou:

\noindent \textit{\textquotedblleft{}Tu nunca mais vai ter um amigo.\textquotedblright{}}

\noindent \textit{\textquotedblleft{}Cala a boca.\textquotedblright{}}

\noindent \textit{\textquotedblleft{}É verdade. Tu é velho. Tarado. Ninguém quer ser amigo de um velho tarado.\textquotedblright{}}

\noindent \textit{\textquotedblleft{}Cala a boca.\textquotedblright{}}

\noindent \textit{\textquotedblleft{}Severino era o único. E ele morreu. Porque todo mundo que chega perto de ti morre.\textquotedblright{}}

\noindent \textit{\textquotedblleft{}Cala a boca!\textquotedblright{}}

\noindent Eu bati a mão na parede. A mão sangrou.

\noindent \textit{\textquotedblleft{}Não foi minha culpa.\textquotedblright{}}

\noindent \textit{\textquotedblleft{}Foi.\textquotedblright{}}

\noindent \textit{\textquotedblleft{}Não foi.\textquotedblright{}}

\noindent \textit{\textquotedblleft{}Severino pegou pneumonia no inverno passado. Sabe o que causou a pneumonia? O frio. Sabe onde o frio entrava? Pelas frestas da parede. A parede que separa a cela dele da sua. As frestas são do lado dele. Mas o frio veio da sua cela.\textquotedblright{}}

\noindent \textit{\textquotedblleft{}Mentira.\textquotedblright{}}

\noindent \textit{\textquotedblleft{}Verdade. Tu respirou o ar frio, o ar passou pela parede, chegou nele. Tu matou ele sem querer. Mas matou.\textquotedblright{}}

\noindent Eu não respondi.

\noindent Fiquei encolhido no canto, a mão sangrando, o coração apertado.

\noindent Ela estava mentindo.
Ela sempre mentia.

\noindent Mas, naquela noite, eu senti o frio entrando pelas frestas.

\noindent E pensei: \textit{\textquotedblleft{}E se for verdade?\textquotedblright{}}

\noindent Mas ali, naquela noite, eu fiz uma escolha. Não a escolha de acreditar nela ou em mim — a escolha de lembrar o Severino como ele era: o único amigo que tive. A criatura podia mentir sobre o frio, sobre a pneumonia, sobre a parede. Eu podia escolher não ouvir. Podia escolher o que guardar dele. E eu guardei a batida na parede. As 5 batidas. A linguagem que a gente criou junto. Isso, ela não podia roubar.

\vspace{0.5em}\noindent\rule{\textwidth}{0.5pt}\vspace{0.5em}

\noindent Nunca mais fiz amizade com ninguém.

\noindent Não valia a pena.
Todo mundo que chegava perto de mim --- a criatura encontrava um jeito de destruir.

\noindent A solidão não era minha escolha.
Era o veneno que ela espalhava em volta de mim.

\noindent Para que eu dependesse só dela.
Para que eu não tivesse mais ninguém.

\noindent Ela queria ser minha única companhia.

\noindent E conseguiu.

\vspace{0.5em}
\noindent\textit{\textbf{↪ Rotas:} \rota{MEM-19} \rotasep \rota{DOC-04}}